% Additional translation from three visibly crossed-out French manuscript blocks.
% Source: archive_20261121/03_sources/EGA5_francais.pdf, physical pp. 126, 136, 140.
% Source SHA-256: 24cf8204184b58b18a46e5f63ecc36d6e791b3c8c96a09daac93de96ebd8e830
% This is a draft-layer appendix, not a restoration of the deleted assertions
% as adopted theorems.  See integration/late_deleted_source_blocks.json.
\begingroup
\def\newscope{V-D}
\def\OO{\mathcal O}
\def\LL{\mathcal L}
\def\NN{\mathcal N}
\def\PP{\mathbf P}
\def\Div{\operatorname{Div}}
\def\Sym{\operatorname{Sym}}
\def\pr{\operatorname{pr}}
\def\unread{\ensuremath{[\![}\mbox{未辨}\ensuremath{]\!]}}

\head{第三篇删划草稿补译（非定稿）}
\NewNote{本章译录法文物理第126、136、140页中被整块大叉划去的三段旧稿。它们作为稿层材料保留；划销这一事实本身不说明作者为什么删去，也不等于作者采用了下列译校修订。明确错误的处理写在对应脚注中；仍不能可靠辨认的字迹以\unread{}标出。第140页旧推论按所印弱假设不成立，正文明确其不能采用的状态，并将原断言完整保留在译校脚注。}

\egasection{法文物理第126页：除子族的旧引入}
\source{法文扫描物理第126页，手写页码133；从“设$X$在$S$上平坦且局部有限表示”至“至迟在第五章”的大叉划销块}

设$X$在$S$上平坦且局部有限表示，$T$为一个$S$-预概形。所谓$X/S$上由$T$参数化的除子族（除子默认为有效 Cartier 除子），是指$X\times_S T$的一个闭子预概形，它是相对于$T$的 Cartier 除子，也就是说，在其每一点处，它在$X\times_S T$中是相对于$T$的余维一的相对完全交。\cmt{V18-D126-01}{原打字句写“由$S$参数化”，而同句环境、$X\times_S T$和后面的函子值均使用$T$，故订为$T$。“$X/S$上”和“相对于$T$的 Cartier 除子，也就是说”等为原页可辨的手写插入；被覆打的旧词不据猜测恢复。}

等价地，它在$T$上平坦且局部有限表示，并在$X_T$的每条纤维
\[
 X_{T,t}=X\times_S\operatorname{Spec}\kappa(t)
\]
上诱导一个“有效 Cartier 除子”，即在其每一点处都为$X_{T,t}$中的余维一完全交的闭子预概形。\cmt{V18-D126-02}{“局部有限表示”及“在$X_t$中”是手写补入。这里的完全交指上述闭浸入的相对 Cartier 条件，不另断言任意$X/S$都是完全交态射。原式把取剩余域谱简写为与$k(t)$作纤维积，正文明写谱。}

所有这类除子族组成的集合记为$\Div_{X/S}(T)$。令$T$变动，便得到一个以集合范畴为值的反变函子，记为$\Div_{X/S}$。我们将在第五章看到，当$f:X\to S$为射影态射时，这个函子是可表出的。

当$X=P=\PP(E)$是由有限型局部自由模$E$定义的射影丛，且$E$各点的秩至少为$2$时，只用第三章的结果，就可以证明$\Div_{P/S}$由下列不交和预概形表出：\cmt{V18-D126-03}{扫描式的主干为$\PP(\Sym^n(E))$，末端有一个细小笔迹，现不能据此可靠核定是否为对偶记号。在本文的一维商约定下，参数化次数$n$有效除子的模必须取$(\Sym^n E)^\vee$，见第三篇2.2—2.4；正文据此明写对偶，不将小笔迹擅自认作清楚的原符号。删稿未明写秩至少为$2$。若将所述不交和用于全体有效除子函子，须限定这一情形：秩$1$时$P\simeq S$，相对有效除子只有空除子，而不同次数的参数分支会重复给出它。次数$n$附近的短小手写补记仍有未辨字，不据此补造秩条件。}
\[
 \coprod_{n\geq 0}\PP\bigl((\Sym^n E)^\vee\bigr).
\]
此和的各部分对应于相对次数为$n$的除子。\cmt{V18-D126-04}{有效除子的次数取$n\geq0$；$n=0$给出空除子。原页在$n$附近另有短小手写，残读为\unread{}；本处不把它当作已核定的句子。}
这个结果对后文并非不可缺少，只须把它作为评注指出即可；我们或许会把它收入第三章的某处，至迟则收入第五章。

\egasection{法文物理第136页：关联可逆模的旧论证}
\source{法文扫描物理第136页，手写页码143；下半页自“命题2(b)向我们表明”至“与除子线性系相关联的可逆模”的大叉划销块}

命题2(b)向我们表明：对除子线性系$D$加上一个补充假设（即$D$\unread{}），就能典范地把$X$上的一个可逆模$\LL$与它联系起来；该模是\cmt{V18-D136-01}{本段保留删稿所用的旧编号“2(b)”。括号内关于$D$的条件在右边缘缺损，随后再次出现的条件也未能可靠辨定；不把它猜成$D\in\mathcal D$或$D\ne\varnothing$。第一式的拉回按上下文写作$f^*$。这里$P$是线性系所给映射的靶，参数丛为$P^\vee$。}
\[
 \LL_U=f^*\OO_P(1),\qquad U=X\setminus Z,
\]
的典范延拓。我们还可以如下给出另一个解释，它不依赖于假设$D$\unread{}。

按构造（或几乎按构造），$\LL_U\otimes_{\OO_S}\OO_{P^\vee}(1)$典范同构于$D$所定义的、$X\times_S P^\vee$上的可逆模$\NN$到$U\times_S P^\vee$的限制$\NN'$。由第1号，这在与该同构相容的唯一同构意义下确定了$\LL_U$；而且（由于$U$在$X$中稠密），存在$X$上的一个可逆模$\LL$，也在相容的唯一同构意义下确定，从而延拓$\LL_U$，并有同构\cmt{V18-D136-02}{张量积中省写了到纤维积的拉回，下面用两投影明写。原文两次说“至唯一同构”；这里的唯一性须连同已经给出的分解同构理解。原删稿把全局存在性缩写为“因为$U$在$X$中稠密”；稠密性本身不足以延拓任意可逆层，此处还须使用第1号对全局可逆模$\NN$的射影丛分解，在稠密开集上确定相对于参数丛的扭转次数。该短语不构成一般的稠密开集延拓定理。}
\[
 \pr_1^*\LL\otimes\pr_2^*\OO_{P^\vee}(1)\simeq\NN.
\]
例如，可以由下式给出$\LL$：\cmt{V18-D136-03}{原删稿的式子是$\LL=\pr_{1*}(\NN)$，少了参数侧的负扭转。依前面的分解，应先张量$\pr_2^*\OO_{P^\vee}(-1)$，再沿$\pr_1$推前。在参数射影丛正秩、即各纤维非空的既定情形，$\pr_{1*}\OO=\OO_X$，投影公式给出正文的恢复式；它不适用于以空参数丛来恢复$\LL$。未扭转的推前通常还带有参数向量丛因子。}
\[
 \LL\simeq\pr_{1*}\!\left(
 \NN\otimes\pr_2^*\OO_{P^\vee}(-1)\right).
\]
我们称它为与除子线性系$D$相关联的、$X$上的可逆模。

\NewNote{本页上半、现行推论3.5证明附近还存在约三行密集覆打的旧句，另有被横线划去的“并且，由于”及余维式。其余维式现只能残读为$\operatorname{codim}(X-U,X)\ \unread\ 2$；比较符号、覆打前的完整假设和相连句子未复原。本章不把这些痕迹补成一段完整证明。}

\egasection{法文物理第140页：旧推论5.3及未竟解释}
\source{法文扫描物理第140页，手写页码147；中部自“以下结果”至“两截面的情形”的整块划销文字}

以下结果在技术上比5.2简单一些，因此应当放在它之前：

\egas{5.3}{原删划推论的处理}
\textbf{按删稿所印的弱假设，这条断言不能采用。} 原来要断言的是：由射影丛$P^\vee$参数化的除子线性系所给的态射
\[
 P^\vee\longrightarrow\Div_{X/S}
\]
为单态射。

本版可用的加强条件来自后面的命题4.3(a)，而不是作者在这一旧稿中已作的修订：在原有忠实平坦、局部有限表示条件之外，要求$f:X\to S$有限表示，并且对每个$s\in S$及$\kappa(s)$的代数闭包$k$，几何纤维$X_k$中存在一个关联素循环，其约化闭包$T$满足$k\xrightarrow{\sim}H^0(T,\OO_T)$。在这些条件下，以上单态射结论成立。\cmt{V18-D140-01}{旧推论的完整陈述译为：“推论5.3。设$f:X\to S$为忠实平坦且局部有限表示的态射。那么，对$X/S$上任一由$S$上射影丛$P^\vee$参数化的除子线性系，相应态射$P^\vee\to\Div_{X/S}$为单态射。”它的弱假设不够：取$S=\operatorname{Spec}k$及$X=\PP^1_k\setminus\{a,b\}$，其中$a\ne b$为两个$k$-有理点；限制点除子的标准线性系时，$a$、$b$都给出$X$上的空除子。这正是4.3所列的反例。此处用4.3(a)说明可用的加强方式，未据划销标记推断作者删改的原因。}

\noindent\textbf{原删稿随后所作的解释（译）。}
这就是说，若两个态射$S'\to P^\vee$与上述态射有相同的复合，它们就相等。作基变换$S'\to S$，便归结为$P^\vee$在$S$上两个截面的情形，\cmt{V18-D140-02}{划销块在页上止于这个逗号；随后另起保留段落，并没有给出这条旧推论的完整证明。左侧斜向边注可辨若干零散字词及如$X_{\overline k}$的记号，但句法与它们所对应的插入位置仍未核定，记为\unread{}，不拼造为作者的新假设。}

\par\endgroup
